\subsection{The Kummer cover extends across the toric cusp}
\label{cusp-cover-section}

This section determines exactly what is true in the assertion that the
degree-$N^2$ cover of the regular fibres extends by a toric power map.  The
answer is positive for the completed \emph{local} cusp family, but the source
action has to be changed, and two different degrees must not be conflated.

Retain the cusp data
\[
 B_0=\begin{pmatrix}0&1\\-1&0\end{pmatrix},\qquad
 Z(s)=sB_0+C(t),\qquad t=e^{2\pi i s},
\]
where $C$ is holomorphic on a disc $\Delta_t$ about zero.  Put
\[
 c_{\bar\lambda}(t)
 :=\exp\!\bigl(2\pi iC(t)\bar\lambda\bigr)
 \in(\mathbb C^*)^2,
 \qquad \bar\lambda\in\bar\Lambda\simeq\mathbb Z^2.
\]
Let $\Sigma$ be the fan in
$N'=\mathbb Z^2\oplus\mathbb Z$ consisting of the cones over the cells of
the translation-invariant $A_2$ triangulation at height one, and write
$Y=Y_\Sigma$.  On the dense torus, with coordinates $(x_1,x_2,t)$, the
target action is
\begin{equation}
 \Psi_{\bar\lambda}(x,t)
 =\bigl(c_{\bar\lambda}(t)t^{B_0\bar\lambda}x,t\bigr).
 \label{cusp-cover-target-action}
\end{equation}
Its quotient is the completed cusp family
$f_0\colon N_0=Y/\bar\Lambda\to\Delta_t$.

Fix an integer $N\geq1$.  Let $v$ be a coordinate on a second disc and
make the base change
\begin{equation}
 \rho_N\colon\Delta_v\longrightarrow\Delta_t,\qquad
 \rho_N(v)=v^N,\qquad t=v^N.
 \label{cusp-cover-base-change}
\end{equation}
On a copy $Y^{(N)}$ of the same toric variety, with dense-torus coordinates
$(\xi_1,\xi_2,v)$, define
\begin{equation}
 C_N(v):=\frac1N C(v^N),\qquad
 c^{(N)}_{\bar\lambda}(v)
 :=\exp\!\bigl(2\pi iC_N(v)\bar\lambda\bigr)
 \label{cusp-cover-source-factor}
\end{equation}
and
\begin{equation}
 \Psi^{(N)}_{\bar\lambda}(\xi,v)
 :=\bigl(c^{(N)}_{\bar\lambda}(v)
          v^{B_0\bar\lambda}\xi,v\bigr).
 \label{cusp-cover-source-action}
\end{equation}

\begin{theorem}[Exact Kummer extension through the cusp]
\label{cusp-cover-extension-theorem}
\mbox{\Verified}\ 
After shrinking the $v$-disc, the following statements hold.

\begin{enumerate}
\item The maps in \eqref{cusp-cover-source-action} extend to a free,
properly discontinuous action on $Y^{(N)}$.  The quotient
\[
 f_N\colon N_0^{\langle N\rangle}
 :=Y^{(N)}/\bar\Lambda\longrightarrow\Delta_v
\]
is a proper holomorphic map from a smooth threefold.  Its central fibre
$W_N$ is reduced, irreducible, and nonnormal, and its normalization is a
toric del Pezzo surface of degree six.

\item Multiplication by $N$ on the cocharacter lattice defines a finite
flat toric morphism
\begin{equation}
 P_N\colon Y^{(N)}\longrightarrow Y,\qquad
 P_N(\xi_1,\xi_2,v)=(\xi_1^N,\xi_2^N,v^N),
 \label{cusp-cover-toric-power-map}
\end{equation}
of degree $N^3$.  It is equivariant for
\eqref{cusp-cover-source-action} and
\eqref{cusp-cover-target-action}, hence descends to a finite flat map
\begin{equation}
 \kappa_N\colon N_0^{\langle N\rangle}\longrightarrow N_0,
 \qquad f_0\circ\kappa_N=\rho_N\circ f_N,
 \label{cusp-cover-completed-map}
\end{equation}
of degree $N^3$.

\item The map $\kappa_N$ is etale over $v\ne0$, and
\begin{equation}
 \kappa_N^*W=NW_N,
 \qquad
 \operatorname{Ram}(\kappa_N)=(N-1)W_N.
 \label{cusp-cover-direct-ramification}
\end{equation}
In particular, the branch support of the map to the original cusp family
is the whole central fibre.

\item Let
\begin{equation}
 N_0[N]:=N_0\times_{\Delta_t}\Delta_v,
 \qquad t=v^N.
 \label{cusp-cover-base-changed-target}
\end{equation}
Then $N_0[N]$ is normal and $\kappa_N$ factors uniquely as
\begin{equation}
 N_0^{\langle N\rangle}
 \xrightarrow{\ Q_N\ }N_0[N]
 \longrightarrow N_0.
 \label{cusp-cover-relative-factorization}
\end{equation}
The map $Q_N$ has generic degree $N^2$.  For $N>1$ it is quasi-etale but
not flat: it is etale over the smooth locus of $N_0[N]$, and its non-etale
and nonflat loci are exactly the inverse images of the double and triple
strata of the base-changed central fibre.  Thus $Q_N$ is not a
normalization map.

\item On the reduced central fibres, the induced map
$W_N\to W$ has generic degree $N^2$.  It lifts to a commutative diagram
\begin{equation}
\begin{CD}
 E_N\simeq\mathrm{dP}_6 @>{[N]}>> E\simeq\mathrm{dP}_6\\
 @V{\nu_N}VV @VV{\nu}V\\
 W_N @>>> W,
\end{CD}
 \label{cusp-cover-normalization-diagram}
\end{equation}
where $[N]$ is toric multiplication by $N$, of degree $N^2$, ramified
with index $N$ along each of the six toric boundary curves.  The diagram
respects the paired conductor curves and the two triple-point classes.

\item Over $v\ne0$, $Q_N$ is exactly the cover defined by
\[
 \mathbb Z\langle\widehat\gamma,\widehat u,
 N\widehat w,N\widehat\delta\rangle
 \subset
 \mathbb Z\langle\widehat\gamma,\widehat u,
 \widehat w,\widehat\delta\rangle.
\]
Its deck group is $(\mathbb Z/N\mathbb Z)^2$.  The pulled-back cusp
monodromy is $M_0^N$, and in the displayed cover basis it is $M_0$:
\begin{equation}
 S_N^{-1}M_0^NS_N=M_0,\qquad
 S_N=\operatorname{diag}(1,1,N,N).
 \label{cusp-cover-monodromy-conjugacy}
\end{equation}
\end{enumerate}
\end{theorem}

\begin{proof}
We prove the assertions in the order in which the constructions occur.

\medskip\noindent
\emph{The source action.}
For $\bar\lambda\in\bar\Lambda$, the integral linear map
\[
 \varphi_{\bar\lambda}(y,y_3)
 =(y+y_3B_0\bar\lambda,y_3)
\]
translates the height-one slice by the integral vector
$B_0\bar\lambda$.  It preserves the $A_2$ triangulation and hence induces
a toric automorphism $\Phi_{\bar\lambda}$ of $Y^{(N)}$.  On the dense
torus this automorphism is
\[
 \Phi_{\bar\lambda}(\xi,v)
 =(v^{B_0\bar\lambda}\xi,v).
\]
Multiplication by the nowhere-zero holomorphic function
$c^{(N)}_{\bar\lambda}(v)$ extends across every toric chart.  Since
$C_N(v)\bar\lambda$ is additive in $\bar\lambda$, the resulting maps
satisfy
$\Psi^{(N)}_{\bar\lambda}\Psi^{(N)}_{\bar\mu}
=\Psi^{(N)}_{\bar\lambda+\bar\mu}$.

For completeness, freeness and proper discontinuity can be checked
directly.  Put
\[
 R_N(v):=-2\pi\operatorname{Im}C_N(v).
\]
It is bounded at $v=0$.  At a dense-torus point fixed by
$\bar\lambda$, logarithmic absolute values give
\[
 \log|v|B_0\bar\lambda+R_N(v)\bar\lambda=0.
\]
Choose the disc so small that
$|\log|v||>2\|R_N\|_\infty$.  Because $B_0$ is orthogonal, the last
equality is impossible for $\bar\lambda\ne0$.  On the central fibre,
$\Psi^{(N)}_{\bar\lambda}$ translates every bounded cell by the nonzero
vector $B_0\bar\lambda$, so it fixes no toric orbit.

For a dense-torus point set
\[
 y(\xi,v)=\frac{(\log|\xi_1|,\log|\xi_2|)}{\log|v|},
 \qquad
 B_{N,v}=B_0+\frac{R_N(v)}{\log|v|}.
\]
Then
\[
 y(\Psi^{(N)}_{\bar\lambda}(\xi,v))
 =y(\xi,v)+B_{N,v}\bar\lambda,
 \qquad
 |B_{N,v}\bar\lambda|\ge\tfrac12|\bar\lambda|.
\]
The logarithmic-position regions associated with any two fixed toric
charts are bounded.  Hence only finitely many $\bar\lambda$ can translate
one of them to meet the other.  Local finiteness of the fan reduces the
same assertion for compact sets to finitely many chart pairs, proving
proper discontinuity.  Finally, reducing
$B_{N,v}^{-1}y(\xi,v)$ modulo $\mathbb Z^2$ places every orbit meeting a
closed subdisc in finitely many closed unit toric charts.  Their union is
compact.  This proves properness of $f_N$.  The central-fibre and
normalization assertions follow exactly from the height-one fan: locally
$v$ is the product of the coordinates belonging to the rays through a
cell, every component occurs with multiplicity one, the lattice action is
transitive on vertices, and the star of a vertex is the complete smooth
hexagonal fan.

\medskip\noindent
\emph{The toric map and equivariance.}
Multiplication by $N$ on $N'$ sends every cone $\sigma\in\Sigma$ onto
itself as a real cone.  It therefore defines the toric morphism $P_N$.
For a maximal cone, let $z_0,z_1,z_2$ be the target characters dual to its
three primitive height-one ray generators, and let $w_0,w_1,w_2$ be the
corresponding source characters.  The map on convergent local rings is
\begin{equation}
 \mathbb C\{z_0,z_1,z_2\}\longrightarrow
\mathbb C\{w_0,w_1,w_2\},\qquad z_i\longmapsto w_i^N.
 \label{cusp-cover-maximal-chart-map}
\end{equation}
The target-chart preimage is the corresponding source chart, and the
right-hand ring is free of rank $N^3$ over the left-hand ring, with basis
\[
 \{w_0^{a_0}w_1^{a_1}w_2^{a_2}:0\le a_i<N\}.
\]
Thus $P_N$ is finite flat of degree $N^3$ on every maximal chart, hence
globally.

The choice \eqref{cusp-cover-source-factor} is forced by equivariance.
Indeed, writing powers coordinatewise,
\begin{align}
 P_N\Psi^{(N)}_{\bar\lambda}(\xi,v)
 &=\bigl((c^{(N)}_{\bar\lambda}(v))^N
          v^{NB_0\bar\lambda}\xi^N,v^N\bigr)\\
 &=\bigl(c_{\bar\lambda}(v^N)
          (v^N)^{B_0\bar\lambda}\xi^N,v^N\bigr)\\
 &=\Psi_{\bar\lambda}P_N(\xi,v).
 \label{cusp-cover-equivariance-computation}
\end{align}
Consequently $P_N$ descends to \eqref{cusp-cover-completed-map}.  To see
finiteness and flatness after the infinite-lattice quotient without any
implicit descent theorem, take an evenly covered open set in $N_0$ and
one of its lifts $V\subset Y$.  Equivariance gives a unique representative
of every source orbit above that open set whose image lies in $V$; two
such representatives cannot differ by a nontrivial lattice translate.
The descended map is therefore locally identified with
$P_N^{-1}(V)\to V$.  It is finite flat of degree $N^3$ by
\eqref{cusp-cover-maximal-chart-map}.

In the same chart the base characters are
\[
 t=z_0z_1z_2,\qquad v=w_0w_1w_2,\qquad P_N^*t=v^N.
\]
The Jacobian determinant of \eqref{cusp-cover-maximal-chart-map} is
\[
 N^3(w_0w_1w_2)^{N-1}=N^3v^{N-1}.
\]
This proves \eqref{cusp-cover-direct-ramification}; the local quotient
maps are biholomorphisms, so the calculation is unchanged downstairs.

\medskip\noindent
\emph{The relative degree-$N^2$ factor.}
The equality
\((f_0\circ\kappa_N)([\xi,v])=v^N=(\rho_N\circ f_N)([\xi,v])\)
gives the factorization
\eqref{cusp-cover-relative-factorization}.  At a triple point the exact
local-ring map is
\begin{equation}
 \begin{gathered}
 A_N:=
 \frac{\mathbb C\{z_0,z_1,z_2,v\}}
      {(z_0z_1z_2-v^N)},
 \qquad
 B:=\mathbb C\{w_0,w_1,w_2\},\\
 A_N\longrightarrow B,\qquad
 z_i\longmapsto w_i^N,\qquad
 v\longmapsto w_0w_1w_2.
 \end{gathered}
 \label{cusp-cover-basechange-local-ring}
\end{equation}
Let
\[
 G_N=\{(\varepsilon_0,\varepsilon_1,\varepsilon_2)\in\mu_N^3:
       \varepsilon_0\varepsilon_1\varepsilon_2=1\}
 \simeq\mu_N^2
\]
act diagonally on $B$.  A monomial
$w_0^{a_0}w_1^{a_1}w_2^{a_2}$ is $G_N$-invariant if and only if
\[
 a_0\equiv a_1\equiv a_2\pmod N.
\]
Writing their common residue as $r$, with $0\le r<N$, expresses that
monomial uniquely in the form
\[
 (w_0w_1w_2)^r
 (w_0^N)^{b_0}(w_1^N)^{b_1}(w_2^N)^{b_2}.
\]
The same coefficientwise argument applies to convergent power series.
Hence
\begin{equation}
 A_N=B^{G_N}.
 \label{cusp-cover-invariant-ring}
\end{equation}
It follows that $A_N$ is normal, that the map has generic degree
$|G_N|=N^2$, and that it is a finite quotient rather than a normalization.
The corresponding face-chart calculations prove normality along the
double and smooth strata as well, so $N_0[N]$ is normal everywhere.

The stabilizer in $G_N$ is trivial when none or exactly one of the $w_i$
vanishes, is isomorphic to $\mu_N$ when exactly two vanish, and is all of
$G_N$ when all three vanish.  Thus $Q_N$ is etale in codimension one and
has no ramification divisor, but it is not etale above the double and
triple strata.  It is also nonflat at every such point when $N>1$.  At a
triple point the closed fibre of \eqref{cusp-cover-basechange-local-ring}
has length
\begin{align}
 \dim_{\mathbb C}
 \frac{\mathbb C\{w_0,w_1,w_2\}}
      {(w_0^N,w_1^N,w_2^N,w_0w_1w_2)}
 &=N^3-(N-1)^3\\
 &=3N^2-3N+1,
 \label{cusp-cover-triple-fibre-length}
\end{align}
which exceeds the generic rank $N^2$.  At a point of a double stratum,
localizing at the nonzero third coordinate and choosing its holomorphic
$N$th root reduces the ramified factor to
\[
 \frac{\mathbb C\{z_0,z_1,v\}}{(z_0z_1-v^N)}
\longrightarrow\mathbb C\{w_0,w_1\},\qquad
 z_i\longmapsto w_i^N,\quad v\longmapsto w_0w_1.
\]
Its local generic rank is $N$, whereas its closed fibre has length
\begin{equation}
 \dim_{\mathbb C}
 \frac{\mathbb C\{w_0,w_1\}}
      {(w_0^N,w_1^N,w_0w_1)}=2N-1>N.
 \label{cusp-cover-double-fibre-length}
\end{equation}
The remaining transverse nonzero torus coordinate is etale locally.
Equations \eqref{cusp-cover-triple-fibre-length} and
\eqref{cusp-cover-double-fibre-length}, together with etaleness elsewhere,
identify the nonflat locus exactly.

\medskip\noindent
\emph{The regular fibres and monodromy.}
Choose a logarithm $r$ with $v=e^{2\pi ir}$.  The target period matrix at
$t=v^N$ is
\[
 Z(Nr)=NrB_0+C(v^N),
\]
whereas the source residual action is determined by
\[
 Z_N(r)=rB_0+\frac1NC(v^N)=\frac1N Z(Nr).
\]
The source fibre is therefore
\[
 \frac{\mathbb C^2}{\mathbb Z^2+Z_N(r)\mathbb Z^2},
\]
and multiplication by $N$ on $\mathbb C^2$ maps it to the target fibre.
Equivalently, in the unscaled additive coordinate it is the identity map
\begin{equation}
 \frac{\mathbb C^2}
      {N\mathbb Z^2+Z(Nr)\mathbb Z^2}
 \longrightarrow
 \frac{\mathbb C^2}
      {\mathbb Z^2+Z(Nr)\mathbb Z^2}.
 \label{cusp-cover-additive-fibre-map}
\end{equation}
The imaginary part of $Z(Nr)$ is invertible on a regular fibre, so its
period lattice meets $\mathbb Z^2$ trivially in the required real sense;
the index in \eqref{cusp-cover-additive-fibre-map} is $N^2$.  Its deck
group is $(\mathbb Z/N\mathbb Z)^2$.  In the ordered lattice basis
$(\widehat\gamma,\widehat u,\widehat w,\widehat\delta)$, this is exactly
the sublattice obtained by $S_N$.  Since $M_0=I+R$ with $R^2=0$ and
$R(a,b,c,d)=(0,0,b,-a)$, one has
\[
 M_0^N=I+NR,\qquad S_N^{-1}M_0^NS_N=M_0,
\]
proving \eqref{cusp-cover-monodromy-conjugacy}.

\medskip\noindent
\emph{The central normalization and conductor.}
Multiplication by $N$ fixes every ray of $\Sigma$.  It therefore maps a
chosen source component of the toric central fibre to the corresponding
target component.  The star of that ray is the hexagonal fan, so the
restriction is the toric power map
$[N]\colon\mathrm{dP}_6\to\mathrm{dP}_6$.  On its dense two-torus it has
degree $N^2$; in each toric boundary chart it is an $N$th-power map in the
normal coordinate, hence has ramification index $N$ along all six boundary
curves.  Equation \eqref{cusp-cover-equivariance-computation} says that
the source and target lattice identifications commute with this map.
Consequently opposite boundary curves, their conductor identifications,
and the two translation-orbits of triangle vertices descend compatibly.
This proves \eqref{cusp-cover-normalization-diagram} and completes the
proof.
\end{proof}

\begin{remark}[Why the unmodified action does not work]
\label{cusp-cover-naive-action-failure}
If one reuses the target action with $v$ in place of $t$ at the source,
then the two sides of equivariance have multiplicative factors
\[
 c_{\bar\lambda}(v)^N v^{NB_0\bar\lambda}
 \quad\hbox{and}\quad
 c_{\bar\lambda}(v^N)v^{NB_0\bar\lambda},
\]
respectively.  In general
$c_{\bar\lambda}(v)^N\ne c_{\bar\lambda}(v^N)$.  Thus the bare formula
$(\xi_1,\xi_2,v)\mapsto(\xi_1^N,\xi_2^N,v^N)$ extends on the ambient toric
variety but does not descend through the manuscript's quotient unless the
source factor \eqref{cusp-cover-source-factor} is inserted.  The explicit
logarithm $2\pi iC(v^N)\bar\lambda/N$ makes this root holomorphic; no
unjustified choice of an $N$th root is involved.
\end{remark}

\begin{remark}[Exact scope]
\label{cusp-cover-global-scope}
The theorem constructs a finite morphism of completed families over the
local cusp discs.  It is a degree-$N^3$ map to the original family and a
degree-$N^2$ map only after the target is base changed.  It is neither a
degree-$N^2$ self-map over the unchanged disc nor the normalization of the
base change.

Nothing here supplies a cover over the entire base $B$.  In fact the
displayed sublattice is not preserved by the original full monodromy when
$N>1$: for example, the printed matrix $A_1$ sends
$\widehat u$ to a vector whose $\widehat w$- and
$\widehat\delta$-coefficients are $-1$ and $1$.  A further global base
cover and compatible extensions through both finite fillings would have
to be constructed separately.  The local Kummer morphism also supplies
no map to a product sphere, smash sphere, $S^{12}$, or $S^{24}$.
\end{remark}
